% GENERATED FILE. Edit canonical lessons, then run npm run generate:books.
% canonical-source-hash: fnv1a64:5ad1f2cb5cbfb2a7
% canonical-lessons: HI-C505-sok, HI-C505-smasan, HI-C505-cita, HI-C505-vidhva, HI-C505-dafnana, HI-R505-review-499-505, HI-R505-review-final-505

\chapter{Mourning}
\label{ch:hi-mourning}
\hlchaptermodality{505}

\begin{chapteropening}
\textbf{By the end of this chapter:} \emph{I can talk about a death in the family.}

Complete the last lesson of chapter 505: I can talk about a death in the family.
\end{chapteropening}

\section[\texorpdfstring{\hi{शोक} (śok)}{शोक (śok)}]{\hi{शोक} (śok) — mourning; grief}

\label{lesson:HI-C505-sok}



\begin{quote}
Before the new one: say the Hindi for to slip from memory, then the Hindi for an inheritance.
\end{quote}

\subsection*{You'll want to know: \hi{शोक}}
\textbf{\hi{शोक}} — \emph{śok} — \textquotedblleft{}mourning; grief\textquotedblright{}.

\begin{quote}
\hi{दादा} \hi{की} \hi{मौत} \hi{पर} \hi{पूरे} \hi{गाँव} \hi{ने} \hi{शोक} \hi{मनाया।}
\end{quote}

\begin{quote}
\emph{dādā kī maut par pūre gā̃v ne śok manāyā.}
\end{quote}

\begin{quote}
The whole village mourned my grandfather's death.
\end{quote}

Masculine.

\begin{cousinweb}
\textbf{\hi{शोक}} (\emph{śok}) is a tatsama, borrowed from Sanskrit \emph{śoka}, \textquotedblleft{}sorrow, anguish, grief\textquotedblright{} (Monier-Williams). \emph{śok manānā} is to observe mourning.
\end{cousinweb}

\begin{grammarlens}[title={What you've built}]
One more word for describing what you lived and felt.
\end{grammarlens}

\subsection*{Your turn}
\begin{itemize}
\raggedright
  \item \emph{Say it:} \emph{śok}
  \item \emph{Say it:} \emph{śok}, once more
  \item \emph{Say it:} \emph{dādā kī maut par pūre gā̃v ne śok manāyā.} — The whole village mourned my grandfather's death.
  \item \emph{Recall:} say the Hindi for to slip from memory, then the Hindi for an inheritance, then say \emph{śok} again
\end{itemize}

\begin{tcolorbox}[breakable,colback=teal!4,colframe=teal!35!black,title={Before you move on}]
What does \hi{शोक} mean? (\textquotedblleft{}Mourning; grief\textquotedblright{}.) Say it once more.
\end{tcolorbox}
\section[\texorpdfstring{\hi{श्मशान} (śmaśān)}{श्मशान (śmaśān)}]{\hi{श्मशान} (śmaśān) — a cremation ground}

\label{lesson:HI-C505-smasan}



\begin{quote}
Before the new one: say the Hindi for an inheritance, then the Hindi for mourning.
\end{quote}

\subsection*{You'll want to know: \hi{श्मशान}}
\textbf{\hi{श्मशान}} — \emph{śmaśān} — \textquotedblleft{}a cremation ground\textquotedblright{}.

\begin{quote}
\hi{परिवार} \hi{के} \hi{लोग} \hi{श्मशान} \hi{तक} \hi{गए।}
\end{quote}

\begin{quote}
\emph{parivār ke log śmaśān tak gae.}
\end{quote}

\begin{quote}
The family went as far as the cremation ground.
\end{quote}

Masculine.

\begin{cousinweb}
\textbf{\hi{श्मशान}} (\emph{śmaśān}) is a tatsama, borrowed from Sanskrit \emph{śmaśāna}, \textquotedblleft{}a place for burning dead bodies\textquotedblright{} (Monier-Williams).
\end{cousinweb}

\begin{grammarlens}[title={What you've built}]
One more word for describing what you lived and felt.
\end{grammarlens}

\subsection*{Your turn}
\begin{itemize}
\raggedright
  \item \emph{Say it:} \emph{śmaśān}
  \item \emph{Say it:} \emph{śmaśān}, once more
  \item \emph{Say it:} \emph{parivār ke log śmaśān tak gae.} — The family went as far as the cremation ground.
  \item \emph{Recall:} say the Hindi for an inheritance, then the Hindi for mourning, then say \emph{śmaśān} again
\end{itemize}

\begin{tcolorbox}[breakable,colback=teal!4,colframe=teal!35!black,title={Before you move on}]
What does \hi{श्मशान} mean? (\textquotedblleft{}A cremation ground\textquotedblright{}.) Say it once more.
\end{tcolorbox}
\section[\texorpdfstring{\hi{चिता} (citā)}{चिता (citā)}]{\hi{चिता} (citā) — a funeral pyre}

\label{lesson:HI-C505-cita}



\begin{quote}
Before the new one: say the Hindi for mourning, then the Hindi for a cremation ground.
\end{quote}

\subsection*{You'll want to know: \hi{चिता}}
\textbf{\hi{चिता}} — \emph{citā} — \textquotedblleft{}a funeral pyre\textquotedblright{}.

\begin{quote}
\hi{बेटे} \hi{ने} \hi{पिता} \hi{की} \hi{चिता} \hi{को} \hi{आग} \hi{दी।}
\end{quote}

\begin{quote}
\emph{beṭe ne pitā kī citā ko āg dī.}
\end{quote}

\begin{quote}
The son lit his father's funeral pyre.
\end{quote}

Feminine.

\begin{cousinweb}
\textbf{\hi{चिता}} (\emph{citā}) is a tatsama, borrowed from Sanskrit \emph{citā}, \textquotedblleft{}a layer, a pile of wood, a funeral pile\textquotedblright{} (Monier-Williams).
\end{cousinweb}

\begin{grammarlens}[title={What you've built}]
One more word for describing what you lived and felt.
\end{grammarlens}

\subsection*{Your turn}
\begin{itemize}
\raggedright
  \item \emph{Say it:} \emph{citā}
  \item \emph{Say it:} \emph{citā}, once more
  \item \emph{Say it:} \emph{beṭe ne pitā kī citā ko āg dī.} — The son lit his father's funeral pyre.
  \item \emph{Recall:} say the Hindi for mourning, then the Hindi for a cremation ground, then say \emph{citā} again
\end{itemize}

\begin{tcolorbox}[breakable,colback=teal!4,colframe=teal!35!black,title={Before you move on}]
What does \hi{चिता} mean? (\textquotedblleft{}A funeral pyre\textquotedblright{}.) Say it once more.
\end{tcolorbox}
\section[\texorpdfstring{\hi{विधवा} (vidhvā)}{विधवा (vidhvā)}]{\hi{विधवा} (vidhvā) — a widow}

\label{lesson:HI-C505-vidhva}



\begin{quote}
Before the new one: say the Hindi for a cremation ground, then the Hindi for a funeral pyre.
\end{quote}

\subsection*{You'll want to know: \hi{विधवा}}
\textbf{\hi{विधवा}} — \emph{vidhvā} — \textquotedblleft{}a widow\textquotedblright{}.

\begin{quote}
\hi{मीरा} \hi{की} \hi{दादी} \hi{तीस} \hi{साल} \hi{से} \hi{विधवा} \hi{हैं।}
\end{quote}

\begin{quote}
\emph{Mīrā kī dādī tīs sāl se vidhvā haiṁ.}
\end{quote}

\begin{quote}
Mira's grandmother has been a widow for thirty years.
\end{quote}

Feminine.

\begin{cousinweb}
\textbf{\hi{विधवा}} (\emph{vidhvā}) is a tatsama, borrowed from Sanskrit \emph{vidhavā}, \textquotedblleft{}a widow\textquotedblright{} (Monier-Williams).
\end{cousinweb}

\begin{grammarlens}[title={What you've built}]
One more word for describing what you lived and felt.
\end{grammarlens}

\subsection*{Your turn}
\begin{itemize}
\raggedright
  \item \emph{Say it:} \emph{vidhvā}
  \item \emph{Say it:} \emph{vidhvā}, once more
  \item \emph{Say it:} \emph{Mīrā kī dādī tīs sāl se vidhvā haiṁ.} — Mira's grandmother has been a widow for thirty years.
  \item \emph{Recall:} say the Hindi for a cremation ground, then the Hindi for a funeral pyre, then say \emph{vidhvā} again
\end{itemize}

\begin{tcolorbox}[breakable,colback=teal!4,colframe=teal!35!black,title={Before you move on}]
What does \hi{विधवा} mean? (\textquotedblleft{}A widow\textquotedblright{}.) Say it once more.
\end{tcolorbox}
\section[\texorpdfstring{\hi{दफ़नाना} (dafnānā)}{दफ़नाना (dafnānā)}]{\hi{दफ़नाना} (dafnānā) — to bury (the dead)}

\label{lesson:HI-C505-dafnana}



\begin{quote}
Before the new one: say the Hindi for a funeral pyre, then the Hindi for a widow.
\end{quote}

\subsection*{You'll want to know: \hi{दफ़नाना}}
\textbf{\hi{दफ़नाना}} — \emph{dafnānā} — \textquotedblleft{}to bury (the dead)\textquotedblright{}.

\begin{quote}
\hi{चिड़िया} \hi{मर} \hi{गई}, \hi{तो} \hi{बच्चों} \hi{ने} \hi{उसे} \hi{बाग़} \hi{में} \hi{दफ़नाया।}
\end{quote}

\begin{quote}
\emph{ciṛiyā mar gaī, to baccoṁ ne use bāġ meṁ dafnāyā.}
\end{quote}

\begin{quote}
The bird died, so the children buried it in the garden.
\end{quote}

Past, with \hi{ने} (he or she …): \hi{उसने} \hi{चिड़िया} \hi{को} \hi{दफ़नाया}, \emph{usne ciṛiyā ko dafnāyā}.

\begin{cousinweb}
\textbf{\hi{दफ़नाना}} (\emph{dafnānā}) is \emph{dafn}, \textquotedblleft{}burial\textquotedblright{}, made into a verb with \emph{-ānā}; \emph{dafn} is usually traced to Arabic \emph{dafn}, \textquotedblleft{}burial\textquotedblright{}, which reached Hindi through Persian.
\end{cousinweb}

\begin{grammarlens}[title={What you've built}]
One more word for describing what you lived and felt.
\end{grammarlens}

\subsection*{Your turn}
\begin{itemize}
\raggedright
  \item \emph{Say it:} \emph{dafnānā}
  \item \emph{Say it:} \emph{dafnānā}, once more
  \item \emph{Say it:} \emph{ciṛiyā mar gaī, to baccoṁ ne use bāġ meṁ dafnāyā.} — The bird died, so the children buried it in the garden.
  \item \emph{Recall:} say the Hindi for a funeral pyre, then the Hindi for a widow, then say \emph{dafnānā} again
\end{itemize}

\begin{tcolorbox}[breakable,colback=teal!4,colframe=teal!35!black,title={Before you move on}]
What does \hi{दफ़नाना} mean? (\textquotedblleft{}To bury (the dead)\textquotedblright{}.) Say it once more.
\end{tcolorbox}
\section[(dialogue)]{Looking back: from Growing Up to Mourning}

\label{lesson:HI-R505-review-499-505}



\begin{quote}
Say the Hindi for an adolescent and of age.
\end{quote}

\subsection*{Your turn}
Say these aloud:

\begin{itemize}
\raggedright
  \item an adolescent, of age, mischief, the prime of youth, to whimper
  \item a bridegroom, henna, an engagement (to marry), the groom's wedding procession, a farewell
  \item a bond, to marry, a divorce, to keep up, a reconciliation
  \item a grandson (a son's son), a grandson (a daughter's son), a nephew (a sister's son), to rock, only (child)
  \item a household, the head (of a family or village), an elder, a generation, to nurture
  \item memory (the power to remember), a walking stick, a pension, to slip from memory, an inheritance
  \item mourning, a cremation ground, a funeral pyre, a widow, to bury (the dead)
\end{itemize}

\begin{tcolorbox}[breakable,colback=teal!4,colframe=teal!35!black,title={Before you move on}]
An adolescent? (\textbf{\hi{किशोर}}, \emph{kiśor}.) Of age? (\textbf{\hi{बालिग़}}, \emph{bāliġ}.) Mischief? (\textbf{\hi{शरारत}}, \emph{śarārat}.) The prime of youth? (\textbf{\hi{यौवन}}, \emph{yauvan}.) To whimper? (\textbf{\hi{बिसूरना}}, \emph{bisūrnā}.) A bridegroom? (\textbf{\hi{दूल्हा}}, \emph{dūlhā}.) Henna? (\textbf{\hi{मेहंदी}}, \emph{mehndī}.) An engagement (to marry)? (\textbf{\hi{सगाई}}, \emph{sagāī}.) The groom's wedding procession? (\textbf{\hi{बारात}}, \emph{bārāt}.) A farewell? (\textbf{\hi{विदाई}}, \emph{vidāī}.) A bond? (\textbf{\hi{नाता}}, \emph{nātā}.) To marry? (\textbf{\hi{ब्याहना}}, \emph{byāhnā}.) A divorce? (\textbf{\hi{तलाक़}}, \emph{talāq}.) To keep up? (\textbf{\hi{निभाना}}, \emph{nibhānā}.) A reconciliation? (\textbf{\hi{सुलह}}, \emph{sulah}.) A grandson (a son's son)? (\textbf{\hi{पोता}}, \emph{potā}.) A grandson (a daughter's son)? (\textbf{\hi{नाती}}, \emph{nātī}.) A nephew (a sister's son)? (\textbf{\hi{भांजा}}, \emph{bhāñjā}.) To rock? (\textbf{\hi{झुलाना}}, \emph{jhulānā}.) Only (child)? (\textbf{\hi{इकलौता}}, \emph{iklautā}.) A household? (\textbf{\hi{गृहस्थी}}, \emph{gṛhasthī}.) The head (of a family or village)? (\textbf{\hi{मुखिया}}, \emph{mukhiyā}.) An elder? (\textbf{\hi{बुज़ुर्ग}}, \emph{buzurg}.) A generation? (\textbf{\hi{पीढ़ी}}, \emph{pīṛhī}.) To nurture? (\textbf{\hi{पोसना}}, \emph{posnā}.) Memory (the power to remember)? (\textbf{\hi{याददाश्त}}, \emph{yāddāśt}.) A walking stick? (\textbf{\hi{छड़ी}}, \emph{chaṛī}.) A pension? (\textbf{\hi{पेंशन}}, \emph{peṁśan}.) To slip from memory? (\textbf{\hi{बिसरना}}, \emph{bisarnā}.) An inheritance? (\textbf{\hi{विरासत}}, \emph{virāsat}.) Mourning? (\textbf{\hi{शोक}}, \emph{śok}.) A cremation ground? (\textbf{\hi{श्मशान}}, \emph{śmaśān}.) A funeral pyre? (\textbf{\hi{चिता}}, \emph{citā}.) A widow? (\textbf{\hi{विधवा}}, \emph{vidhvā}.) To bury (the dead)? (\textbf{\hi{दफ़नाना}}, \emph{dafnānā}.)
\end{tcolorbox}
\section[(dialogue)]{Mourning — once more}

\label{lesson:HI-R505-review-final-505}



\begin{quote}
Say the Hindi for mourning and a cremation ground.
\end{quote}

\subsection*{Your turn}
\begin{itemize}
\raggedright
  \item \emph{Say it:} mourning, a cremation ground, a funeral pyre, a widow, to bury (the dead)
\end{itemize}

\begin{tcolorbox}[breakable,colback=teal!4,colframe=teal!35!black,title={Before you move on}]
Mourning? (\textbf{\hi{शोक}}, \emph{śok}.) A cremation ground? (\textbf{\hi{श्मशान}}, \emph{śmaśān}.) A funeral pyre? (\textbf{\hi{चिता}}, \emph{citā}.) A widow? (\textbf{\hi{विधवा}}, \emph{vidhvā}.) To bury (the dead)? (\textbf{\hi{दफ़नाना}}, \emph{dafnānā}.)
\end{tcolorbox}
